\section{The character estimate over a cyclotomic field of arbitrary allowed level}

% BEGIN T82K1
\paragraph{T82K1. Restriction to the full cyclotomic field and the torsion obstruction.}
Use the data in $\mathrm{Cyc}(N,K,L)$ and set
$s=\max(1,v_p(N))$, $F=\mathbb Q_p(\mu_{\mathrm{lcm}(p,N)})$.
The coefficient field contains $F$, and $K\subset F$. The
prime-to-$p$ roots generate an unramified field $F_0$, and
$F=F_0(\mu_{p^s})$. At $p=2$ the hypothesis $4\mid N$ gives
$s\ge2$. Write $G_s=G_F$, $H=G_{F_0(\mu_{p^\infty})}$.
The image of $G_s/H$ under $\chi$ is $1+p^s\mathbb Z_p$.
Choose $\gamma$ lifting its topological generator $u=1+p^s$.
At these levels $v_p(u^j-1)=s+v_p(j)$ for each nonzero integer $j$.

In the first test of $\mathrm{TBAt}$ the integral functional to the
trivial line has an integral lift of $1$. Its kernel lattice $T$ is
saturated and the differences of the lift give the connecting cocycle
$f$. The rational kernel is crystalline on restriction to $G_s$,
and its jumps lie among the allowed ambient jumps by strictness. The
comparison W79H1--W79H5 therefore gives its Wach hypothesis with
the stated height bounds. Restriction of a crystalline filtration
extends scalars at each embedding and hence preserves its jumps.
Let $e<t$, $M=e+1$, $A=A_{\inf}$, $A_M=A/(\xi^M)$ and
$B_M=B_{\rm dR}^+/\mathrm{Fil}^{M}=A_M[1/p]$.
For $m=2s+n_0(p,2h)+1$ put $c=\mathrm{cLL}(p,m,m-s,J(T),0,M)$.
The invariant lift in filtration zero used in K81C1 gives a
rational primitive for $f\otimes1$ on $G_s$. Clearing its
denominators and applying W81T6 gives an integral primitive for
$p^cf\otimes1$ on the same group.

Choose any $\mathbb Z_p$-coordinate on the cyclotomic target
lattice and lift the tested map linearly on $T$. Write $c_f(g)$ for
the corresponding coordinate evaluated on $f(g)$. K81C2's
contraction of the integral primitive yields
$p^c c_f(g)-D_g b\in A_M$, where $D_g b=\chi(g)^{-e}g b-b$.
The failure of equivariance of the coordinate lift is integral
on $T$, which suffices in the contraction. The reduction of $c_f$
is a torsion cocycle. Local duality identifies the dual of
$H^2(F,\mathbb Z_p(-e))$ with $(\mathbb Q_p/\mathbb Z_p(e+1))^{G_s}$.
Testing $\gamma$ kills this group by $p^d$, $d=s+v_p(e+1)$.
Exactness of the cohomology sequence therefore gives a rational
cocycle lifting $p^d c_f$ modulo integral scalars, after a
coboundary adjustment. Multiplying by $p^c$ gives a rational
cocycle $Z$ with
\[
 Z(g)\equiv p^{c+d}c_f(g)\pmod{\mathbb Z_p},
 \qquad Z(g)-D_g(p^db)\in A_M .
\]
% END T82K1

% BEGIN T82K2
\paragraph{T82K2. An upper bound for the logarithmic coefficient.}
Every rational scalar $\chi^{-e}$-cocycle on $G_s$ has, in $B_M$,
a unique logarithmic coefficient $\alpha\in F$, meaning that its
difference from $g\mapsto\alpha\log\chi(g)t^e$ is a coboundary.
Here $t=\log[\epsilon]$ denotes the period in this section's
formulas for $B_M$. If the cocycle is integral modulo a coboundary
in $A_M$, then
\[
                   |\alpha|\le p^{e(s+1)+v_p(e!)+s+1+e/(p-1)} .
\]
Indeed the integral trace-$p$ averaging of A79L1 and K81C3 on $H$
uses the same infinite tower for every finite level. After multiplying
a bounded cocycle by $p$ and removing its $H$-coboundary, use the
integral normalized-trace splitting onto $F$ in A79L3 at $m=s$.
On its complement $u^j\gamma-1$ has inverse of norm at most
$p^{1/(p-1)}$ for every integer $j$. On $F$ it has inverse of norm
$p^{s+v_p(j)}$ for $j\ne0$. Thus a unit-ball cocycle on a nonzero
grade $C(j)$ has a unique primitive of norm at most $p^{s+1+v_p(j)}$.
For $j=0$ subtract the logarithmic cocycle determined by the
$F$-component at $\gamma$. Division by $p\log u$ bounds its
coefficient by $p^{s+1}$. A logarithmic coboundary has zero
coefficient, by normalized trace.

The successive grades in $B_M(-e)$ have weights $-e,\ldots,0$.
The preceding primitives remove the lower grades, and the weight-zero
calculation gives existence and uniqueness on the final grade. For the
integral bound, the $t^i$-residue of a primitive in degree $i$ is
multiplied by $(\zeta_p-1)^i$ when expressed in the integral
$\xi^i$-coordinate, exactly as in K81C3. This scalar is fixed by
$G_s$. If the difference is integral, clearing the bound just proved
allows one to lift that residue integrally and increase the filtration
by one. The accumulated exponent in degrees $i<e$ is
$e(s+1)+v_p(e!)$. In the top degree apply the weight-zero bound and
divide by $(\zeta_p-1)^e$. Its absolute value is $p^{-e/(p-1)}$,
giving the assertion.
% END T82K2

% BEGIN T82K3
\paragraph{T82K3. A Frobenius equation and an integral primitive at the dyadic prime.}
Use $E=[\epsilon^{1/p^s}]$, $\mu=[\epsilon]-1$,
$\pi_j=[\epsilon^{1/p^j}]-1$, $\xi=\mu/\pi_1$ and $U=u^{-e}$.
Let $\sigma$ denote Frobenius on $F_0$. For a finite set of pairs
$(c_i,m_i)$, $c_i\in\mathcal O_{F_0}$, $m_i>0$ prime to $p$, put
$Q=\sum_i c_i E^{m_i}$. Suppose
$(1-\sigma)d_0=\sigma^s(\sum_i c_i)$ with $d_0$ integral.
An integral solution of $(1-\varphi)y=Q$ is
\[
 y=\sum_{k=0}^{s-1}\varphi^k Q+d_0+
       \sum_{n\ge0}\sum_i\sigma^{s+n}(c_i)((1+\mu)^{m_i p^n}-1).
\]
The coefficients in each fixed positive degree of the tail tend
to zero $p$-adically, while high degrees carry the corresponding
powers of $\mu$. Thus the series converges in the complete $(p,\mu)$
topology, and the Frobenius difference telescopes. In $B_M$ write
$y=\sum_{j=0}^e w_jt^j$. Since
$\varphi^k E=\zeta_{p^{s-k}}\exp(t/p^{s-k})$ for $k<s$,
\[
 w_j=\frac{
 \sum_{k=0}^{s-1}p^{kj}\sum_i\sigma^k(c_i)m_i^j\zeta_{p^{s-k}}^{m_i}
              +p^{sj}d_j}{j!p^{sj}},
 \qquad
 d_j=(1-p^j\sigma)^{-1}\sigma^s\!\left(\sum_i c_i m_i^j\right)
 \quad(j>0),
\]
with the given $d_0$ for $j=0$.

There is an element $b_0\in B_M$ satisfying
$D_h b_0\in\mathbb Z_p$ for $h\in H$ and
$D_\gamma b_0+\varphi^{-1}y\in\mathbb Z_p$, where the inverse
Frobenius is first applied in $A$. To see this, use the integral
recursion of K81C4 with $X=\mu$ and $v=((1+X)^u-1)/X$. We have
$\gamma E=E(1+X)$ and $v\equiv1\bmod(p,X^2)$, since
$p^s\ge3$. On a monomial $E^m$, the first nonzero coefficient
of the operator
$P\mapsto U(1+X)^mP(Xv)-v^{k+1}P(X)$ on $X^i$ modulo $p$ is
$m$ in degree $i+1$. Thus successive coefficients solve the
recursion for every unit exponent $m$. In particular, starting with
$Q_0=Q$, one obtains $H$-fixed $P_k,Q_k$ in $A$ satisfying
\[
 D_\gamma(P_k/\mu^{k+1})=
        Q_k/\mu^k-p Q_{k+1}/\mu^{k+1}.
\]
Choose integral $\beta_k,\eta_k$ by
$(\varphi-\xi^{k+1})\beta_k=P_k$ and
$(\varphi-\xi^k)\eta_k=Q_k$, with $\eta_0=-y$. These equations
are solvable by lifting roots of $z^p-a z-b$ in the algebraically
closed tilt, one Witt digit at a time. Dividing by $\pi_1$ to the
corresponding powers changes the Frobenius differences to differences
with denominator $\mu$. Consequently
\[
 D_\gamma(\beta_k/\pi_1^{k+1})-\eta_k/\pi_1^k
                       +p\eta_{k+1}/\pi_1^{k+1}\in\mathbb Z_p,
 \qquad D_h(\beta_k/\pi_1^{k+1})\in\mathbb Z_p .
\]
These identities live first in $W(C^\flat)$, whose Frobenius-fixed
ring is $\mathbb Z_p$.

Apply $\varphi^{-1}$ there before passing the identities to $B_M$.
It fixes the scalar errors, commutes with $G_s$, and changes the
denominator to $\pi_2$. Clearing denominators in $A$ justifies
passage to $B_M$, where their residues are nonzero. Write $D=p(p-1)$.
The shifted cyclotomic polynomial gives
$\xi=\Phi_{p^2}(1+\pi_2)=\pi_2^D+p(1+\pi_2 w)$ with $w\in A$.
Modulo $(\pi_2^a,\xi^M)$ the second factor is a unit and
$(\xi-\pi_2^D)^{\lceil a/D\rceil+M-1}=0$ by the binomial theorem.
It follows that
$p^{\lceil a/D\rceil+M-1}\pi_2^{-a}\in A_M$ inside $B_M$.
Here $D\ge2$ even for $p=2$. Therefore
\[
 b_0=\sum_{k\ge0}p^k\varphi^{-1}(\beta_k)/\pi_2^{k+1}
\]
converges in the $p$-adically complete $A_M[1/p]$, and the final
$\eta$-term tends to zero. Telescoping proves the two assertions for
$b_0$ since the scalar lattice is closed.
% END T82K3

% BEGIN T82K4
\paragraph{T82K4. Removing the extra Frobenius term and realizing the coefficient.}
From the integral equation,
$\varphi^{-1}y=y+Q'$ with $Q'=\varphi^{-1}Q$. This extra term has
an $H$-fixed twisted primitive in $B_M$. Indeed its coefficient of
$t^j$ is a sum of rational multiples of
$\sigma^{-1}(c_i) m_i^j\zeta_{p^{s+1}}^{m_i}$ in
$F'=F_0(\mu_{p^{s+1}})$. Each has trace zero over $F$, since
the $m_i$ are units and the conjugates differ by all the $p$-th
roots of unity. The restriction $\gamma'$ of $\gamma$ generates
$\operatorname{Gal}(F'/F)$. On its trace kernel the map
$u^{j-e}\gamma'-1$ is invertible. For $j\ne e$ its $p$-fold
geometric sum gives an inverse up to $u^{p(j-e)}-1\ne0$.
For $j=e$ its kernel is the intersection of $F$ and the trace
kernel, which is zero in characteristic zero. Invert these maps
coefficientwise and multiply by $t^j$ to obtain $b_1$ with
$D_\gamma b_1=Q'$ and $D_h b_1=0$. Thus $D_\gamma(b_0+b_1)+y$
is an integral scalar.

Adding $\sum_{j<e}w_jt^j/(u^{j-e}-1)$ to $b_0+b_1$ gives $b_2$
for which $z(g)=D_g b_2+(w_e/\log u)\log\chi(g)t^e$ belongs
to $\mathbb Z_p$ on $H$ and $\gamma$. It is a continuous cocycle
with values in $B_M(-e)$. The inverse image of the closed stable
scalar lattice under it is a closed subgroup, hence contains all
of $G_s$. The resulting integral scalar cocycle has logarithmic
coefficient $w_e/\log u$.
% END T82K4

% BEGIN T82K5
\paragraph{T82K5. The integral coefficient lattice over all primitive branches.}
Let $a$ run through the unit residue classes modulo $p^s$.
For integral $c_a,d\in\mathcal O_{F_0}$ satisfying
$(1-p^e\sigma)d=\sigma^s\sum_a c_a$, consider the additive map
\[
 T_{s,e}(c,d)=\sum_{k=0}^{s-1}p^{ke}\sum_a\sigma^k(c_a)
                  \zeta_{p^{s-k}}^a+p^{se}d .
\]
Its image contains $p^{(s-1)(e+1)}\mathcal O_F$. For $s=1$ the
integral basis $\zeta_p^a$, $1\le a<p$, has sum $-1$.
For $x=\sum b_a\zeta_p^a$ solve
$(1-p^{e+1}\sigma)d=\sigma\sum b_a$ by a geometric series
and put $c_a=b_a+p^ed$. This gives $T_{1,e}=x$, including the
algebraic base case $p=2$.

For $s>1$, given one unit class $a$ at level $s$, assign coefficient
$pb$ to it and subtract $b$ at each of the $p$ lifts of its reduction
at level $s-1$. The sum of the top roots on the lifts is zero and
all lower root contributions cancel. Thus the image contains
$p b\zeta_{p^s}^a$ for all integral $b\in F_0$. Next, from any
level-$(s-1)$ data $c'_a,d'$ repeat $\sigma^{-1}c'_a$ at each
lift and take $d=pd'$. These satisfy the new constraint and give
$p^{e+1}T_{s-1,e}(c',d')$. The powers of $\zeta_{p^s}$ generate
the integer ring over $\mathcal O_{F_0}$ by the Eisenstein
uniformizer $\zeta_{p^s}-1$. Each power with exponent divisible
by $p$ is already at level $s-1$. The two constructions give the
assertion by induction, using their additivity and integral
$\mathbb Z_p$-scaling.

For $e=0$ choose the pairs $(c_a,a)$ in T82K3, using positive
unit representatives and $d_0=d$. For $e>0$ take
$(-c_a/a^e,a)$ and $(c_a/a^e,ua)$. Then the zeroth moment
vanishes and $d_0=0$ works. Each $e$-th moment gains the factor
$u^e-1$ while the roots in the numerator are unchanged. The
coefficients in T82K4 are respectively $T_{s,0}/\log u$ and
$(u^e-1)T_{s,e}/(\log(u)e!p^{se})$. In consequence the image of
integral cocycles contains $p^\beta\mathcal O_F$, where
\[
 \beta=-1\quad(e=0),\qquad
 \beta=s-1-e-v_p((e-1)!)\quad(e>0).
\]
% END T82K5

% BEGIN T82K6
\paragraph{T82K6. The original invariant and character constants over the tested field.}
Put $a=(e+2)s+v_p(e)+\lfloor e/(p-1)\rfloor+1$, taking
$v_p(0)=0$ here. The norm bound in T82K2 gives
$p^a\alpha\in p^\beta\mathcal O_F$ by T82K5. For $e>0$ the sum
of the norm exponent and $\beta$ equals
$(e+2)s+v_p(e)+e/(p-1)$, and for $e=0$ the sum is $s$.
Subtract the integral cocycle with that coefficient from $p^a Z$.
The resulting rational cocycle has logarithmic coefficient zero.

The rational cocycle space has dimension $[F:\mathbb Q_p]+1$ by
local duality, the Euler formula and the dimension of the space of
coboundaries (K81C6). The coefficient map is onto $F$ by T82K5
and rational scaling. Its kernel is therefore one-dimensional. For
$e>0$ it contains the nonzero coboundary of $1$. For $e=0$ it
contains the unramified degree homomorphism, whose crystalline
primitive solves a residue Frobenius difference equation in
unramified Witt vectors as in K81C6. Thus every zero-coefficient
cocycle is crystalline-finite. Combining the coordinates supplies
a finite rational lift of $p^{c+d+a}\phi(f)$ over $F$.

Corestrict to $K$ and divide by the unit part of $[F:K]$.
Corestriction commutes with reduction and with the crystalline
period map, and its composite with restriction on a class from $K$
is the degree. Hence $p^{c+d+a+g}\phi(f)$ is finite over $K$,
where $g=v_p([F:K])$. The $p$-part of $[F:\mathbb Q_p]$ comes
from the ramified degree $p^{s-1}$ and the residue degree, which is
at most the prime-to-$p$ part of $N$ (and is one if that part is one).
Consequently $g\le\lfloor\log_p N\rfloor+1$.

Set $\Lambda=s+\lfloor\log_p(2h)\rfloor+\lfloor\log_p N\rfloor+1$,
with the source convention $\log_p0=0$ for integer logarithms.
If all jumps lie in $U\subset[-h,h]$ of size $\le n$, then
$m\le2s+\lfloor\log_p(2h)\rfloor+2$. On each of the $e+1$
grades, every distinct jump $b\ne i$ costs
$m+v_p(b-i)\le2\Lambda$. Thus
$c\le2\Lambda+3(e+1)(n+1)\Lambda$. Here $e<t\le h$.
Using $v_p(e)\le e$, $v_p(e+1)\le e+1$, and $g\le\Lambda$,
we have $c+d+a+g\le40(t+1)(n+1)\Lambda=c_B^L$.
For the other clauses of $\mathrm{TBAt}$, the positive dual
twist has finite cyclic invariants, since $\chi(G_K)$ contains
$u$. The calculation of M81B3 now gives finite equal to the
unramified classes at $e=0$, finite equal to coboundaries at
$e>0$, and the normalized torsion-length bound
$([K:\mathbb Q_p]+1)c_B^L$ at this exponent. This proves
$\mathrm{TBK}$ for the original test.

For $\mathrm{InvariantsKilledOverCyclotomic}$ restrict an invariant
class of $V/T$ to $G_s$. Its representative tensored with $1$ has
integral differences in $T\otimes\mathcal O_C$. W81T3--W81T5 on
the one-grade window give integrality modulo a rational invariant at
exponent $\mathrm{cLL}(p,m,m-s,J,0,1)$. The rational invariant
space is zero since weight zero occurs at no embedding, by the
graded basis in W81T4. Integrality then implies integrality in the
original lattice, by taking coordinates in a $\mathbb Z_p$-basis.
The above one-grade estimate gives
$2\Lambda+3(n+1)\Lambda\le20(n+1)\Lambda=c_0^L$, also for $h=0$.
This proves the invariant assertion for the same $K,L,N$.
% END T82K6
